\section{The dimension formula}
\label{section-dimension-formula}

\noindent
Recall the definitions of catenary (Definition \ref{definition-catenary})
and universally catenary (Definition \ref{definition-universally-catenary}).

\begin{lemma}
\label{lemma-dimension-formula}
Let $R \to S$ be a ring map.
Let $\mathfrak q$ be a prime of $S$ lying over the prime $\mathfrak p$ of $R$.
Assume that
\begin{enumerate}
\item $R$ is Noetherian,
\item $R \to S$ is of finite type,
\item $R$, $S$ are domains, and
\item $R \subset S$.
\end{enumerate}
Then we have
$$
\text{height}(\mathfrak q)
\leq
\text{height}(\mathfrak p) + \text{trdeg}_R(S)
- \text{trdeg}_{\kappa(\mathfrak p)} \kappa(\mathfrak q)
$$
with equality if $R$ is universally catenary.
\end{lemma}

\begin{proof}
Suppose that $R \subset S' \subset S$, where $S'$ is a finitely generated
$R$-subalgebra of $S$. In this case set $\mathfrak q' = S' \cap \mathfrak q$.
The lemma for the ring maps $R \to S'$ and $S' \to S$ implies the
lemma for $R \to S$ by additivity of transcendence degree in towers
of fields (Fields, Lemma \ref{fields-lemma-transcendence-degree-tower}).
Hence we can use induction on the number of generators
of $S$ over $R$ and reduce to the case where $S$ is generated by
one element over $R$.

\medskip\noindent
Case I: $S = R[x]$ is a polynomial algebra over $R$.
In this case we have $\text{trdeg}_R(S) = 1$.
Also $R \to S$ is flat and hence
$$
\dim(S_{\mathfrak q}) =
\dim(R_{\mathfrak p}) + \dim(S_{\mathfrak q}/\mathfrak pS_{\mathfrak q})
$$
see Lemma \ref{lemma-dimension-base-fibre-equals-total}.
Let $\mathfrak r = \mathfrak pS$. Then
$\text{trdeg}_{\kappa(\mathfrak p)} \kappa(\mathfrak q) = 1$
is equivalent to $\mathfrak q = \mathfrak r$, and implies that
$\dim(S_{\mathfrak q}/\mathfrak pS_{\mathfrak q}) = 0$.
In the same vein $\text{trdeg}_{\kappa(\mathfrak p)} \kappa(\mathfrak q) = 0$
is equivalent to having a strict inclusion
$\mathfrak r \subset \mathfrak q$, which implies that
$\dim(S_{\mathfrak q}/\mathfrak pS_{\mathfrak q}) = 1$.
Thus we are done with case I with equality in every instance.

\medskip\noindent
Case II: $S = R[x]/\mathfrak n$ with $\mathfrak n \not = 0$.
In this case we have $\text{trdeg}_R(S) = 0$.
Denote $\mathfrak q' \subset R[x]$ the prime corresponding to $\mathfrak q$.
Thus we have
$$
S_{\mathfrak q} = (R[x])_{\mathfrak q'}/\mathfrak n(R[x])_{\mathfrak q'}
$$
By the previous case we have
$\dim((R[x])_{\mathfrak q'}) =
\dim(R_{\mathfrak p}) + 1
- \text{trdeg}_{\kappa(\mathfrak p)} \kappa(\mathfrak q)$.
Since $\mathfrak n \not = 0$ we see that the dimension of
$S_{\mathfrak q}$ decreases by at least one, see
Lemma \ref{lemma-one-equation},
which proves the inequality of the lemma.
To see the equality in case $R$ is universally catenary note that
$\mathfrak n \subset R[x]$ is a height one prime as it corresponds
to a nonzero prime in $F[x]$ where $F$ is the fraction field of $R$.
Hence any maximal chain of primes in
$S_\mathfrak q = R[x]_{\mathfrak q'}/\mathfrak nR[x]_{\mathfrak q'}$
corresponds to a maximal chain of primes
with length 1 greater between $\mathfrak q'$ and $(0)$ in $R[x]$.
If $R$ is universally catenary these all have the same length equal to
the height of $\mathfrak q'$. This proves that
$\dim(S_\mathfrak q) = \dim(R[x]_{\mathfrak q'}) - 1$
and this implies equality holds as desired.
\end{proof}

\noindent
The following lemma says that generically finite maps
tend to be quasi-finite in codimension $1$.

\begin{lemma}
\label{lemma-finite-in-codim-1}
Let $A \to B$ be a ring map.
Assume
\begin{enumerate}
\item $A \subset B$ is an extension of domains,
\item the induced extension of fraction fields is finite,
\item $A$ is Noetherian, and
\item $A \to B$ is of finite type.
\end{enumerate}
Let $\mathfrak p \subset A$ be a prime of height $1$.
Then there are at most finitely many primes of $B$
lying over $\mathfrak p$ and they all have height $1$.
\end{lemma}

\begin{proof}
By the dimension formula (Lemma \ref{lemma-dimension-formula})
for any prime $\mathfrak q$ lying over $\mathfrak p$ we have
$$
\dim(B_{\mathfrak q}) \leq
\dim(A_{\mathfrak p}) - \text{trdeg}_{\kappa(\mathfrak p)} \kappa(\mathfrak q).
$$
As the domain $B_\mathfrak q$ has at least $2$ prime ideals we see that
$\dim(B_{\mathfrak q}) \geq 1$. We conclude that
$\dim(B_{\mathfrak q}) = 1$ and that the extension
$\kappa(\mathfrak p) \subset \kappa(\mathfrak q)$ is algebraic.
Hence $\mathfrak q$ defines a closed point of its fibre
$\Spec(B \otimes_A \kappa(\mathfrak p))$, see
Lemma \ref{lemma-finite-residue-extension-closed}.
Since $B \otimes_A \kappa(\mathfrak p)$ is a Noetherian ring
the fibre $\Spec(B \otimes_A \kappa(\mathfrak p))$
is a Noetherian topological space, see
Lemma \ref{lemma-Noetherian-topology}.
A Noetherian topological space consisting of closed points
is finite, see for example
Topology, Lemma \ref{topology-lemma-Noetherian}.
\end{proof}









